\begin{description}
\item[(a)] Drawing scaled diagram is impossible. Rough sketch is accepted.
\item[(b)] There are 10 days and nights per taris year. The obliquity is
  $3^\circ$, which means that the planet's rotation is in the same direction
  as its orbit. Thus, total number of rotations per year is $10 + 1 = 11$.

  \textbf{Note:} The obliquity is positive (similar to the Earth / Mars /
  Jupiter). This means, we have ADD one rotation. Subtracting one rotation by
  assuming opposite rotation (like the Venus) is incorrect.
\item[(c)] By Kepler's third law, $\dfrac{T^2}{R^3} = \text{Constant}$
  \begin{align*}
  \frac{T_{en}^2}{R_{en}^3} &= \frac{T_{ex}^2}{R_{ex}^3} \tag{1}\\
  R_{ex}^3 &= \frac{1.6^2 R_{en}^3}{0.2^2} \tag{2}\\
  R_{ex} &= \sqrt[3]{64}\ \text{endorlengths} \tag{3}\\
         &= 4\ \text{endorlengths} \tag{4}
  \end{align*}
\item[(d)] Using same logic as above
  \begin{align*}
  \frac{T_C^2}{R_C^3} &= \frac{T_T^2}{R_T^3} \tag{5}\\
  T_C^2 &= \frac{9^3 R_T^3 T_T^2}{R_T^3} \tag{6}\\
  T_C &= \sqrt{729}\ \text{tarisyears} \tag{7}\\
      &= 27\ \text{tarisyears} \tag{8}
  \end{align*}
\item[(e)] As Corulus is in Opposition, Sola - Taris - Corulus form straight
  line (in that order). Distance $= 9 - 1 = 8$ tarislengths.
\item[(f)] In the figure, S is Sola, A and B are start of the year positions
  of Taris and Corulus, T and C are their positions after `$n$' days. Angles
  are named from $a$ to $f$. The dashed line is parallel to line SB.
  Triangle(SCT) is used for sine rule as well as answer in the next part.
  Figure is not to the scale.
  % FIGURE NEEDED: two concentric circular orbits around S (Sola) with
  % start-of-year positions A (Taris) and B (Corulus) on a line through S,
  % later positions T and C after n days, angles a-f marked, dashed line
  % through T parallel to SB; 2010 proceedings, solution of long problem 17,
  % p. 50.
  \begin{align*}
  a + b + c &= \pi \tag{9}\\
  b + d + e &= \pi \tag{10}\\
  d &= f + c \tag{11}\\
  f + c &= \frac{2\pi n}{10} \tag{12}\\
  f &= \frac{2\pi n}{270} \tag{13}\\
  \sin b &= 9\sin a \ \text{(By Sine Rule)} \tag{14}\\[6pt]
  e &= \pi - b - d \tag{15}\\
    &= \pi - b - c - f \tag{16}\\
    &= a - f \tag{17}\\[6pt]
  b &= \pi - (a + c) \tag{18}\\
    &= \pi - \left(a + \frac{2\pi n}{10} - \frac{2\pi n}{270}\right) \tag{19}\\
    &= \pi - \left(a + \frac{52\pi n}{270}\right) \tag{20}\\[6pt]
  9\sin a &= \sin\left(\pi - \left(a + \frac{52\pi n}{270}\right)\right)
    \tag{21}\\
    &= \sin\left(a + \frac{52\pi n}{270}\right) \tag{22}\\
    &= \left[\sin a\cos\left(\frac{52\pi n}{270}\right)
       + \cos a\sin\left(\frac{52\pi n}{270}\right)\right] \tag{23}\\
  9 &= \cos\left(\frac{52\pi n}{270}\right)
       + \cot a\sin\left(\frac{52\pi n}{270}\right) \tag{24}\\
  \cot a &= \frac{9 - \cos\left(\frac{52\pi n}{270}\right)}
                 {\sin\left(\frac{52\pi n}{270}\right)} \tag{25}\\
  a &= \tan^{-1}\!\left[\frac{\sin\left(\frac{52\pi n}{270}\right)}
       {9 - \cos\left(\frac{52\pi n}{270}\right)}\right] \tag{26}\\[6pt]
  \lambda &= \pi - e \tag{27}\\
          &= \pi + f - a \tag{28}\\
  \lambda &= \pi + \frac{2\pi n}{270}
     - \tan^{-1}\!\left[\frac{\sin\left(\frac{52\pi n}{270}\right)}
       {9 - \cos\left(\frac{52\pi n}{270}\right)}\right] \tag{29}
  \end{align*}
\item[(g)]
  \begin{align*}
  \text{Area} &= \frac12 \times l(ST) \times l(SC) \times \sin c\\
              &= \frac12 \times 1 \times 9 \times 0.568\\
              &= 2.56
  \end{align*}
  The area is about 3\,(tarislength)$^2$.
\end{description}
